\section{Full proof of the partition principle}
\label{app:generic}

We prove Proposition~\ref{prop:generic} with all quantifiers explicit.
The argument is stated for a compact Riemannian manifold with normalized volume measure.
The same proof applies to any metric probability space that has the required isoperimetric and coarea properties.

For a fixed public random string $r$, the map $S\mapsto A(S;r)$ has finite range because the sample space is finite.
We denote this range by $\cY_r$.
For $y\in\cY_r$, let
\[
 h_{r,y}(\theta)=\Pr_{S\sim P_\theta^m}\{A(S;r)=y\}.
\]
Assumption 1 and data processing imply
\begin{equation}
 |h_{r,y}(\theta)-h_{r,y}(\phi)|
 \le \TV(P_\theta^m,P_\phi^m)
 \le\eta
 \label{eq:app-event}
\end{equation}
whenever $\dist(\theta,\phi)\le R$.
For the categorical family used in the paper, $h_{r,y}$ is continuous.
Continuity also follows locally from the vanishing right side of the KL bound as $\dist(\theta,\phi)\downarrow0$.

Define
\begin{align*}
 L_r&=\int_M\sum_{y\in\cY_r}h_{r,y}(\theta)
               \ell(y,\theta)\,d\mu(\theta),\\
 D_r&=\int_M\left(1-\sum_{y\in\cY_r}h_{r,y}(\theta)^2\right)d\mu(\theta).
\end{align*}
The hypotheses of Proposition~\ref{prop:generic} hold pointwise in $\theta$.
Fubini's theorem therefore gives
\[
 \E_rL_r\le\alpha,
 \qquad
 \E_rD_r\le\rho.
\]
By Markov's inequality,
\[
 \Pr_r\{L_r>3\alpha\}<\frac13,
 \qquad
 \Pr_r\{D_r>3\rho\}<\frac13.
\]
There is consequently a fixed seed $r$ for which both inequalities in \eqref{eq:goodseed} hold.
We fix it below and suppress $r$ from the notation.

\subsection{Canonical cells and coverage}

For every $y\in\cY_r$, define the open canonical cell
\[
 F_y=\{\theta:h_y(\theta)>3/4\}.
\]
Distinct canonical cells are disjoint.
If $\theta\notin\cup_yF_y$, every coordinate of the probability vector $(h_y(\theta))_y$ is at most $3/4$.
Among probability vectors with this constraint, the squared $\ell_2$ norm is maximized by $(3/4,1/4,0,\ldots)$.
Thus
\[
 1-\sum_yh_y(\theta)^2\ge 1-\frac{10}{16}=\frac38.
\]
Since $D_r\le3\rho$,
\begin{equation}
 \mu(N)\le8\rho,
 \qquad
 N:=M\setminus\bigcup_yF_y.
 \label{eq:noncan-volume}
\end{equation}
This proves \eqref{eq:cover}.

For $t\in[0,R]$, put
\[
 G_{y,t}=\{\theta:\dist(\theta,F_y)<t\}
\]
for $t>0$, and take $G_{y,0}=F_y$.
If $\theta\in\overline{G_{y,t}}$, choose a sequence $\phi_n\in F_y$ with
$\dist(\theta,\phi_n)\le t+o(1)$.
Equation~\eqref{eq:app-event} and continuity give
\begin{equation}
 h_y(\theta)\ge 3/4-\eta=\beta>1/2.
 \label{eq:app-beta}
\end{equation}
It follows that the closures of $G_{y,t}$ and $G_{z,t}$ are disjoint when $y\ne z$.

\subsection{The distortion budget}

For any $t\in[0,R]$, disjointness and \eqref{eq:app-beta} give
\begin{align}
 3\alpha
 &\ge \int_M\sum_y h_y(\theta)\ell(y,\theta)d\mu(\theta)\nonumber\\
 &\ge \beta\sum_y\int_{G_{y,t}}\ell(y,\theta)d\mu(\theta).
 \label{eq:app-dist}
\end{align}
Let $\cG_t$ be the indices satisfying
\[
 \frac{1}{\mu(G_{y,t})}
 \int_{G_{y,t}}\ell(y,\theta)d\mu(\theta)
 \le K\alpha.
\]
Empty cells are omitted.
Equation~\eqref{eq:app-dist} implies
\begin{equation}
 \sum_{y\notin\cG_t}\mu(G_{y,t})
 \le\frac{3}{\beta K}.
 \label{eq:app-badmass}
\end{equation}
The thickened cells contain the original canonical cells, so \eqref{eq:noncan-volume} and \eqref{eq:app-badmass} yield
\begin{equation}
 \sum_{y\in\cG_t}\mu(G_{y,t})
 \ge 1-8\rho-\frac{3}{\beta K}.
 \label{eq:app-goodmass}
\end{equation}
For $K\ge32$, $\eta<1/4$, and sufficiently small universal $\rho$, the right side is at least a positive universal constant $c_*$.

Fix $y\in\cG_t$.
Markov's inequality within the cell gives
\[
 \mu\{\theta\in G_{y,t}:\ell(y,\theta)\le2K\alpha\}
 \ge\frac12\mu(G_{y,t}).
\]
Assumption 2 of Proposition~\ref{prop:generic} therefore implies
\begin{equation}
 \mu(G_{y,t})\le2s.
 \label{eq:app-cellsmall}
\end{equation}
Assumption 3, \eqref{eq:app-goodmass}, and \eqref{eq:app-cellsmall} now show that for almost every $t$,
\begin{equation}
 \sum_y\Per_\mu(G_{y,t})
 \ge \lambda\sum_{y\in\cG_t}\mu(G_{y,t})
 \ge c_*\lambda.
 \label{eq:app-per}
\end{equation}

\subsection{Boundary points are noncanonical}

Fix $t\in(0,R]$ and $\theta\in\partial G_{y,t}$.
Then $\theta\notin F_y$, since every point of the open set $F_y$ is interior to $G_{y,t}$.
If $\theta\in F_z$ for some $z\ne y$, then $h_z(\theta)>3/4$ while \eqref{eq:app-beta} gives $h_y(\theta)\ge\beta>1/2$.
This is impossible.
Thus $\theta\in N$.
More generally, every point of the shell $G_{y,R}\setminus F_y$ is on the boundary of $G_{y,t}$ at
$t=\dist(\theta,F_y)$ and belongs to $N$ by the same argument.
The shells for distinct $y$ are disjoint because the closures of $G_{y,R}$ are disjoint.
Consequently,
\begin{equation}
 \sum_y\mu(G_{y,R}\setminus F_y)\le\mu(N)\le8\rho.
 \label{eq:app-shell}
\end{equation}

The distance function to a set has metric gradient one almost everywhere outside its closure.
The coarea formula, first for regular distance levels and then by approximation, gives
\begin{equation}
 \int_0^R\Per_\mu(G_{y,t})dt
 \le\mu(G_{y,R}\setminus F_y).
 \label{eq:app-coarea}
\end{equation}
Summing \eqref{eq:app-coarea}, using \eqref{eq:app-shell}, and integrating \eqref{eq:app-per} gives
\[
 c_*R\lambda
 \le\int_0^R\sum_y\Per_\mu(G_{y,t})dt
 \le8\rho.
\]
This is \eqref{eq:generic} and completes the proof of Proposition~\ref{prop:generic}.

\section{Continuity of the categorical experiment}
\label{app:kl}

We prove Lemma~\ref{lem:kl}.
For brevity, write $c_i=\cos\theta_i$, $c'_i=\cos\phi_i$, and similarly $s_i,s'_i$.
Since $a\le1/2$, every nonzero atom of $p_\phi$ is at least $1/(8d)$.
The inequality $\KL(p\Vert q)\le\chi^2(p,q)$ gives
\begin{align*}
 \KL(p_\theta\Vert p_\phi)
 &\le \sum_{i,j}
 \frac{(p_\theta(i,j)-p_\phi(i,j))^2}{p_\phi(i,j)}\\
 &\le 8d\sum_{i=1}^d
 \frac{2a^2(c_i-c'_i)^2+2a^2(s_i-s'_i)^2}{16d^2}\\
 &=\frac{a^2}{d}\sum_{i=1}^d
 \big((c_i-c'_i)^2+(s_i-s'_i)^2\big).
\end{align*}
The expression in parentheses is the squared chord length between two points on the unit circle.
It is at most $\Delta_i^2$, proving \eqref{eq:kl}.

KL tensorizes over independent samples.
Pinsker's inequality therefore gives
\begin{equation}
 \TV(p_\theta^m,p_\phi^m)
 \le a\sqrt{\frac{m}{2d}}\,dist(\theta,\phi).
 \label{eq:app-pinsker}
\end{equation}
Choose $c<\sqrt2/15$ in the radius of Lemma~\ref{lem:kl}.
Then the right side of \eqref{eq:app-pinsker} is below $1/15$.
The probability of any event, including the event that a fixed-seed estimator returns a specified output, changes by at most this total variation distance.
This proves the lemma.

\section{The small-basin lemma}
\label{app:basin}

We prove Lemma~\ref{lem:basin}.
Fix $q\in\Delta_k$ and define $u_i$ as in the main text.
Let $v(\theta_i)=(\cos\theta_i,\sin\theta_i)$.
If
\[
 L_i(\theta)=\sum_{j=1}^4|q(i,j)-p_\theta(i,j)|,
\]
then the triangle inequality in each signed pair gives
\begin{equation}
 \norm{u_i-v(\theta_i)}_2
 \le\frac{2d}{a}L_i(\theta).
 \label{eq:app-ui}
\end{equation}

We next record the radial projection fact used in \eqref{eq:decode}.
For any $u\in\R^2$, there is an angle $\varphi$ such that, for every $\theta$,
\begin{equation}
 \min\{1,\dist_\T(\theta,\varphi)\}
 \le\pi\norm{u-v(\theta)}_2.
 \label{eq:app-radial}
\end{equation}
If $\norm{u}_2<1/2$, the right side is at least $\pi/2$ and the claim is immediate.
Otherwise choose $v(\varphi)=u/\norm{u}_2$.
Then
\[
 \norm{v(\theta)-v(\varphi)}_2
 \le2\norm{v(\theta)-u}_2.
\]
For an angular separation $x\in[0,\pi]$, the chord length is $2\sin(x/2)\ge2x/\pi$.
Equation~\eqref{eq:app-radial} follows.

Combining \eqref{eq:app-ui} and \eqref{eq:app-radial}, and summing over blocks, yields
\begin{align*}
 \TV(q,p_\theta)
 &=\frac12\sum_{i,j}|q(i,j)-p_\theta(i,j)|\\
 &\ge\frac{a}{4\pi d}\sum_{i=1}^d
 \min\{1,\dist_\T(\theta_i,\varphi_i)\}.
\end{align*}
Mass of $q$ on unused symbols only increases the left side, so the inequality holds for all $q\in\Delta_k$.

For independent uniform angles $\Theta_i$, the circular distance to a fixed $\varphi_i$ is uniform on $[0,\pi]$.
Thus, with $Z_i=\min\{1,\dist_\T(\Theta_i,\varphi_i)\}$,
\[
 \mu_0:=\E Z_i
 =\frac1\pi\left(\int_0^1x\,dx+\int_1^\pi1\,dx\right)
 =1-\frac1{2\pi}.
\]
If $\TV(q,p_\Theta)\le2K\alpha$ and $a=C_0\alpha$, then
\[
 \frac1d\sum_iZ_i\le\frac{8\pi K}{C_0}.
\]
Choose $C_0\ge16\pi K/\mu_0$.
Hoeffding's inequality gives
\[
 \Pr\left\{\frac1d\sum_iZ_i\le\frac{\mu_0}{2}\right\}
 \le\exp\left(-\frac{\mu_0^2d}{2}\right).
\]
Taking $c_0=\mu_0^2/2$ proves \eqref{eq:smallbasin}.

\section{Isoperimetry on the flat torus}
\label{app:torus}

We prove Lemma~\ref{lem:torusiso}.
Let $\gamma_d$ be standard Gaussian measure on $\R^d$ and define
\[
 T(x)_i=2\pi\Phi(x_i)\pmod {2\pi}.
\]
Each coordinate is uniform on the circle, so $T_\#\gamma_d=\mu$.
Moreover,
\[
 \sup_x\left|\frac{d}{dx}2\pi\Phi(x)\right|
 =\sqrt{2\pi}.
\]
Taking shortest circular distances can only reduce distance, hence $T$ is $L=\sqrt{2\pi}$-Lipschitz in the product Euclidean metrics.

Fix measurable $E\subseteq\T^d$ and set $A=T^{-1}(E)$.
Lipschitzness implies
\[
 A_t\subseteq T^{-1}(E_{Lt}).
\]
The Gaussian isoperimetric inequality \citep{borell1975,sudakov1978} states
\[
 \gamma_d(A_t)\ge
 \Phi\big(\Phi^{-1}(\gamma_d(A))+t\big).
\]
Since $\gamma_d(A)=\mu(E)$, we obtain
\[
 \mu(E_{Lt})\ge
 \Phi\big(\Phi^{-1}(\mu(E))+t\big).
\]
Subtracting $\mu(E)$, dividing by $Lt$, and taking a lower limit at zero proves
\[
 \Per_\mu(E)\ge\frac{\phi(\Phi^{-1}(\mu(E)))}{\sqrt{2\pi}}.
\]

For $0<u\le1/2$, standard Gaussian tail bounds imply
\begin{equation}
 \phi(\Phi^{-1}(u))
 \ge c u\sqrt{\log(1/u)}
 \label{eq:app-profile}
\end{equation}
for a universal $c>0$.
If $u\le2e^{-c_0d}$ and $d$ is above a universal constant, then
$\log(1/u)\ge c_0d/2$.
Substitution into \eqref{eq:app-profile} proves \eqref{eq:smalliso}.

\section{Completion of Theorem~\ref{thm:main}}
\label{app:main}

The application of Proposition~\ref{prop:generic} in the main text uses
\[
 R=\min\left\{1,\frac{c_2\sqrt d}{a\sqrt m}\right\}.
\]
Lemma~\ref{lem:kl} verifies sample continuity at this radius after reducing $c_2$ if needed.
When the second term exceeds one, Equation~\eqref{eq:app-pinsker} verifies continuity throughout the unit radius because
$a\sqrt{m/d}<c_2$.

Proposition~\ref{prop:generic} and Lemma~\ref{lem:torusiso} give
\[
 R\sqrt d\le C\rho.
\]
Choose $d_0$ and $\rho_0$ so that $\sqrt d>C\rho$ for $d\ge d_0$ and $\rho\le\rho_0$.
Then $R\ne1$, so
\[
 \frac{c_2d}{a\sqrt m}\le C\rho.
\]
Rearranging and using $a=C_0\alpha$ gives
\[
 m\ge \frac{c_2^2}{C^2C_0^2}
       \frac{d^2}{\alpha^2\rho^2}.
\]
Finally, for $k\ge8$, $d=\lfloor k/4\rfloor\ge k/8$.
Increasing $k_0$ to ensure $d\ge d_0$ proves Theorem~\ref{thm:main}.
